\documentclass[aspectratio=169]{beamer}
\usetheme{Madrid}
\usecolortheme{default}

% Packages
\usepackage[utf8]{inputenc}
\usepackage[portuguese]{babel}
\usepackage{graphicx}
\usepackage{booktabs}

% Title information
\title[Pipeline de Extração de Gabaritos e Visualização em Tempo Real]{Pipeline de OCR Automatizado para Extração de Gabaritos e Visualização em Tempo Real}
\subtitle{Demonstração: Funcionalidades, Fluxos e Otimizações}
\date{\today}

\begin{document}

\begin{frame}
    \titlepage
\end{frame}

\begin{frame}{Objetivo}
    \textbf{O Objetivo:} Automatizar a correção e a digitalização dos cartões-resposta dos alunos com o precisão e velocidade.
    
    \vspace{0.5cm}
    Em vez de corrigir manualmente centenas de provas, o pipeline:
    \begin{itemize}
        \item Lê imagens ou arquivos PDF.
        \item Identifica exatamente o exame, a escola, o professor e os alunos.
        \item Lê as respostas marcadas nas "bolhas" (círculos).
        \item Exporta dados estruturados (CSVs) prontos para análise e dashboards (Google Sheets).
    \end{itemize}
\end{frame}

\begin{frame}{Tipos de Entrada Suportados}
    O sistema é altamente flexível e lida com múltiplos formatos automaticamente:
    
    \vspace{0.3cm}
    \begin{columns}
        \column{0.5\textwidth}
        \textbf{Imagens (PNG, JPG)}
        \begin{itemize}
            \item Fotos tiradas pelos professores ou digitalizações.
            \item A foto deve ser, de preferência, tirada de maneira horizontal e o mais centralizada possível.
            \item Processamento quase instantâneo.
        \end{itemize}
        
        \column{0.5\textwidth}
        \textbf{Documentos (PDFs)}
        \begin{itemize}
            \item Scanners com múltiplas páginas.
            \item O sistema extrai automaticamente as páginas de um PDF em diferentes páginas PNG e as processa.
        \end{itemize}
    \end{columns}
    
    \vspace{0.5cm}
    \begin{block}{Processamento em Lote}
        Você pode apontar o sistema para uma pasta contendo os PDFs e imagens misturadas, e ele processará o lote inteiro de uma só vez.
    \end{block}
\end{frame}

\begin{frame}{Como o Pipeline funciona, de forma rápida}
    \begin{enumerate}
        \item \textbf{Pré-processamento:} A imagem é convertida para tons de cinza e automaticamente endireitada para que tudo fique o mais alinhado possível.
        \item \textbf{Detecção de QR Code:} O sistema escaneia o topo da página buscando um QR Code.
        \item \textbf{Extração da Grade:} Encontramos as linhas verticais e horizontais da tabela de respostas para isolar cada bolha individualmente.
        \item \textbf{Análise das Bolhas:} O sistema analisa as bolhas isoladas (análise de densidade de pixels) para determinar qual opção o aluno marcou.
        \item \textbf{Exportação:} Os resultados são compilados em uma planilha mestre.
    \end{enumerate}
\end{frame}

\begin{frame}{Integração com QR Codes}    
    \vspace{0.5cm}
    \textbf{Informações retiradas}
    \begin{itemize}
        \item O ID do Professor, o ID da Escola e a Turma/Bimestre.
        \item Quais questões exatas estão na página e sua ordem.
        \item Quais alunos estão nesta página específica (IDs dos Alunos).
    \end{itemize}
\end{frame}

\begin{frame}{Como as "bolhas" são detectadas}    
    \vspace{0.3cm}
    \textbf{O método proposto:}
    \begin{itemize}
        \item \textbf{Isolamento:} Recortamos o quadrado aproximado onde a bolha está.
        \item \textbf{Limiarização:} Filtramos a cor de fundo do papel.
        \item \textbf{Checagem de Densidade:} Contamos quantos "pixels escuros" existem dentro da bolha em comparação com uma vazia.
        \item \textbf{Pontuação de Confiança:} Se a densidade for alta (acima de um limiar), temos 100\% de certeza de que foi preenchida.
    \end{itemize}
\end{frame}

\begin{frame}{Otimizando para Velocidade}    
    \vspace{0.3cm}
    Implementei, visando otimizar performance, sem perder a acurácia:
    \begin{enumerate}
        \item Operações com bibliotecas otimizadas (OpenCV e Numpy).
        \item Rejeição de modelos de IA Generativa.
        \item Processamento Paralelo Multi-core (Multiprocessing).
    \end{enumerate}
\end{frame}

\begin{frame}[fragile]{Como Executar: Comandos do Terminal}
    O sistema é controlado através de uma Interface de Linha de Comando (CLI) simples.
    
    \vspace{0.3cm}
    \textbf{Processar uma Única Imagem ou PDF}
    \begin{verbatim}
poetry run python main.py --image "exemplos/prova.pdf"
    \end{verbatim}
    
    \vspace{0.2cm}
    \textbf{Processar um Diretório/Lote Completo (com Multiprocessamento)}
    \begin{verbatim}
poetry run python main.py --batch "exemplos/"
    \end{verbatim}
\end{frame}

\begin{frame}{Métricas Finais de Performance}
    Com todas as otimizações ativas, eis o que o sistema alcança hoje:
    
    \vspace{0.5cm}
    \begin{itemize}
        \item \textbf{Tempo Puro de Extração:} \textbf{0.18 segundos} por página.
        \item \textbf{Velocidade Real em Lote:} Graças ao multiprocessamento, o tempo de espera cai para cerca de \textbf{0.05 segundos} por página.
    \end{itemize}
\end{frame}

\begin{frame}{Exemplos Visuais e Tratamento de Erros}
    \textbf{Transparência e Controle do Processo}
    
    \vspace{0.3cm}
    \begin{itemize}
        \item \textbf{Imagens de Depuração (Debug):} O pipeline pode gerar imagens intermediárias mostrando exatamente onde encontrou o QR Code, a grade da tabela e as caixas delimitadoras de cada "bolha" processada.
        \item \textbf{Antes e Depois:} As imagens permitem visualizar a imagem original (torta ou com sombras) e a versão processada após a correção de perspectiva.
        \item \textbf{Tratamento de Exceções:} Se uma página estiver muito rasgada ou o QR Code for ilegível, o sistema não para. Ele gera um log de erro e separa o arquivo problemático para revisão manual, garantindo que nenhum aluno fique sem nota.
    \end{itemize}
\end{frame}

\begin{frame}{Formato de Saída e Integração}
    \textbf{Dados Estruturados e Prontos para Uso}
    
    \vspace{0.3cm}
    O sistema não entrega apenas notas cruas, mas sim uma visão granular de cada resposta, gerando arquivos \texttt{CSV}. Exemplo de saída:
    
    \vspace{0.2cm}
    \begin{table}[]
        \centering
        \small
        \begin{tabular}{lccccccc}
            \toprule
            \textbf{ID Aluno} & \textbf{Q1} & \textbf{Q2} & \textbf{Q3} & \textbf{Q47} & \textbf{Q5} & \textbf{Q6-A} & \textbf{Q6-B} \\
            \midrule
            A3859 & B & 1 & 2 & 3 & 2 & B & 2 \\
            A3860 & B & 2 & B & 1 & 2 & 2 & 1 \\
            A3861 & B & 1 & 2 & 1 & B & 2 & B \\
            \bottomrule
        \end{tabular}
    \end{table}

    \vspace{0.2cm}
    \begin{itemize}
        \item \textbf{Matrizes de Confiança:} Além da resposta, o pipeline gera arquivos paralelos informando a densidade de preenchimento e o grau de confiança (0-100\%) para cada bolha lida, alertando para possíveis rasuras.
    \end{itemize}
\end{frame}

\begin{frame}{Conclusão}
    \textbf{Objetivos Atingidos}
    Os objetivos combinados foram cumpridos. O processamento terá que ser monitorado, para ajuste de parâmetros,
    mas a implementação está fechada em um ciclo End-to-End (do começo ao fim).
\end{frame}

\end{document}
